\chapter{Concepts: the contracts the DSL enforces}
\label{ch:concepts}

\begin{aside}
Templates can do almost anything. That is their power and their
curse. A library that uses templates freely will, eventually,
produce error messages thousands of lines long where the actual
problem is a single missing function. The C++20 concepts feature
is the cure: a way to say ``this template parameter must satisfy
\emph{this} contract'' and have the compiler enforce it directly.
\maya{} leans on concepts heavily. This chapter walks through the
ones you will encounter and the ones you will write.
\end{aside}

\section{What a concept is, mechanically}

A \emph{concept} is a named compile-time predicate on a type. You
declare it once; you use it as a constraint anywhere a template
parameter appears.

\begin{lstlisting}
template <class T>
concept Addable = requires (T a, T b) {
    { a + b } -> std::convertible_to<T>;
};

template <Addable T>           // T must satisfy Addable
T sum(T a, T b) { return a + b; }
\end{lstlisting}

The \code{requires} clause inside the concept lists expressions
that must compile. \code{a + b} must be a valid expression and
its result must convert to \code{T}. If you call \code{sum(1, 2)},
the compiler checks: does \code{int + int} exist? Yes. Does it
convert to \code{int}? Yes. Substitution proceeds.

If you call \code{sum(std::string\{"a"\}, std::string\{"b"\})}, the
check runs again: does \code{string + string} exist? Yes. Does it
convert to \code{string}? Yes. Proceed.

If you call \code{sum(std::ifstream\{\}, std::ifstream\{\})}, the
check fails: \code{ifstream + ifstream} is not a thing. The
compiler emits a focused error: ``constraint Addable<ifstream>
not satisfied because a + b is ill-formed''.

\subsection*{Before concepts: SFINAE}

Pre-C++20, the same effect was achieved with SFINAE (Substitution
Failure Is Not An Error):

\begin{lstlisting}
template <class T,
          class = std::enable_if_t<std::is_arithmetic_v<T>>>
T sum(T a, T b) { return a + b; }
\end{lstlisting}

This worked, but the error messages were brutal. The constraint
was hidden in a \code{std::enable\_if\_t} that produced no
explanation; failures looked like ``no matching function''.

Concepts fixed three things:

\begin{itemize}[leftmargin=*]
  \item The constraint has a \emph{name} you can see in errors.
  \item The constraint is \emph{first-class}: you can compose
    concepts with \code{\&\&}, \code{||}, \code{!}.
  \item The compiler checks the constraint \emph{before} trying
    to instantiate the function body, so errors point at the
    constraint, not at the failing internals.
\end{itemize}

\begin{funfact}
Concepts were proposed for C++0x and dropped in 2009 because the
design was too complex. A simplified version (Concepts Lite)
appeared in 2014 as TS and finally shipped in C++20. Bjarne
Stroustrup has been writing about concepts since the 1980s; the
wait was personal for him.
\end{funfact}

\section{The concepts in maya's core}

\maya{} defines a handful of concepts in
\file{maya/include/maya/core/concepts.hpp}. Each names a contract
that some other part of the library relies on. Walking them is
the fastest way to understand the library's surface.

\subsection*{Program}

The single most important concept. It states what an application
must provide to be runnable by \code{maya::run}.

\begin{lstlisting}
template <class P>
concept Program = requires {
    typename P::Model;
    typename P::Msg;
    { P::init() }
        -> std::same_as<std::pair<typename P::Model, Cmd<typename P::Msg>>>;
    requires requires (typename P::Model m, typename P::Msg msg) {
        { P::update(m, msg) }
            -> std::same_as<std::pair<typename P::Model, Cmd<typename P::Msg>>>;
        { P::view(m) }      -> std::same_as<Element>;
        { P::subscribe(m) } -> std::same_as<Sub<typename P::Msg>>;
    };
};
\end{lstlisting}

Six requirements:

\begin{enumerate}[leftmargin=*]
  \item \code{P::Model} is a nested type.
  \item \code{P::Msg} is a nested type.
  \item \code{P::init()} is callable and returns a pair of the
    initial model and an initial command.
  \item \code{P::update(model, msg)} returns the same pair shape.
  \item \code{P::view(model)} returns an \code{Element}.
  \item \code{P::subscribe(model)} returns a \code{Sub<Msg>}.
\end{enumerate}

The \code{run} function takes a \code{Program} as a template
parameter:

\begin{lstlisting}
template <Program P>
int run(Config cfg = {});
\end{lstlisting}

If you mistype \code{view} as \code{View}, the compile error names
\code{Program} as the unsatisfied concept and points at the
missing function. No twenty-deep template trace.

\begin{idea}
Concepts in \maya{} are not decorative. They are documentation,
contract, and error-message generator in one. When you read
\code{template <Program P>}, you know exactly what \code{P}
must look like, because the concept is right there.
\end{idea}

\subsection*{ElementLike}

The concept that gates the DSL's children:

\begin{lstlisting}
template <class T>
concept ElementLike = requires (T t) {
    { build(t) } -> std::same_as<Element>;
};
\end{lstlisting}

Any type that has a free function \code{build} returning an
\code{Element} satisfies it. \code{TextNode}, \code{BoxNode},
\code{Element} itself, and the runtime variants all qualify.
This is what lets \code{h(\ldots)} and \code{v(\ldots)} take
heterogeneous children.

\subsection*{Renderable}

Used by widgets that produce their own elements:

\begin{lstlisting}
template <class W>
concept Renderable = requires (const W& w) {
    { w.render() } -> std::same_as<Element>;
};
\end{lstlisting}

Widgets that satisfy \code{Renderable} can be dropped into the
DSL via a \code{widget(\ldots)} helper that calls
\code{w.render()}.

\subsection*{StyleSource}

A type that can supply a \code{Style} on demand. Used by the
theme system:

\begin{lstlisting}
template <class T>
concept StyleSource = requires (const T& t) {
    { t.style() } -> std::same_as<Style>;
};
\end{lstlisting}

A \code{Theme} struct, a styled text node, and a tagged builder
all satisfy this. The renderer asks any \code{StyleSource} for
its style during the merge pass.

\subsection*{Subscribable}

Some pieces of state can be observed for changes:

\begin{lstlisting}
template <class S>
concept Subscribable = requires (S s, std::function<void()> cb) {
    { s.subscribe(cb) } -> std::convertible_to<std::function<void()>>;
};
\end{lstlisting}

The signal system (Chapter~\ref{ch:signals}) uses this to
distinguish reactive sources from plain values. The subscribe
function takes a callback and returns an unsubscribe handle.

\section{Composing concepts}

Concepts compose with the usual logical operators:

\begin{lstlisting}
template <class T>
concept Drawable = ElementLike<T> && Subscribable<T>;

template <class T>
concept Numeric = std::integral<T> || std::floating_point<T>;
\end{lstlisting}

A type satisfies \code{Drawable} if and only if it satisfies both
\code{ElementLike} \emph{and} \code{Subscribable}. The standard
library's \code{std::integral} and \code{std::floating\_point}
are themselves concepts; their union is \code{Numeric}.

You can also negate:

\begin{lstlisting}
template <class T>
concept NotEmpty = !std::same_as<T, std::monostate>;
\end{lstlisting}

\begin{recap}
Concepts are named compile-time predicates. They give templates a
contract the compiler can check up front. \maya{}'s
\code{Program}, \code{ElementLike}, \code{Renderable},
\code{StyleSource}, and \code{Subscribable} concepts describe the
shape of every extensible point in the library.
\end{recap}

\section{Writing your own concept}

Suppose you have a family of components that all expose a
\code{focus()} method returning \code{void} and a \code{focused()}
method returning \code{bool}. You want a function that takes any
such component.

\begin{lstlisting}
template <class C>
concept Focusable = requires (C& c, const C& cc) {
    { c.focus() }      -> std::same_as<void>;
    { cc.focused() }   -> std::same_as<bool>;
};

template <Focusable C>
void cycle_focus(std::vector<C>& cs) {
    for (auto& c : cs)
        if (c.focused()) { c.focus(); return; }
}
\end{lstlisting}

Notice the const-correctness: \code{focused()} is called on a
const reference, so it must be a const method. The concept
catches the mistake of declaring \code{focused()} non-const.

\subsection*{Multiple requires clauses}

You can use multiple \code{requires} clauses for organisation:

\begin{lstlisting}
template <class C>
concept Container = requires (C& c, C const& cc) {
    typename C::value_type;
    typename C::iterator;
    { c.begin() }    -> std::same_as<typename C::iterator>;
    { c.end() }      -> std::same_as<typename C::iterator>;
    { cc.size() }    -> std::convertible_to<std::size_t>;
    { cc.empty() }   -> std::same_as<bool>;
};
\end{lstlisting}

Type aliases (\code{typename C::value\_type}) require the nested
name to exist. Expression checks (\code{c.begin()}) require the
expression to compile, optionally with a return type constraint.

\section{Where concepts appear in maya}

A non-exhaustive list:

\begin{itemize}[leftmargin=*]
  \item \code{template <Program P> int run(Config)} ---
    \file{maya/include/maya/app/app.hpp}.
  \item \code{template <ElementLike\ldots Cs> auto h(Cs\&\& \ldots)}
    --- \file{maya/include/maya/dsl.hpp}.
  \item \code{template <ElementLike\ldots Cs> auto v(Cs\&\& \ldots)}
    --- ditto.
  \item \code{template <Renderable W> auto widget(const W\&)} ---
    DSL helper.
  \item \code{template <class M> Cmd<M> Cmd::task(callable)}, where
    the callable must satisfy a small concept ensuring it returns
    \code{M}.
  \item Platform abstractions: \code{template <Terminal T> class
    Renderer\{\}}, where \code{Terminal} is the concept that the
    POSIX, macOS, and Win32 backends all satisfy.
\end{itemize}

Reading the header, you can see the contract before you see the
implementation. This is what makes \maya{}'s headers approachable.

\begin{funfact}
GCC, Clang, and MSVC implement concepts differently in two
places: the error message format and the order in which concept
checks happen during overload resolution. The standard pins down
the \emph{behaviour}, not the \emph{output}. If you see slightly
different messages on different compilers, that is expected.
\end{funfact}

\section{Pitfalls}

\begin{warning}
\textbf{Concepts can hide failure if you let them.} A
\code{template <class T>} with no constraint and a
\code{requires} clause inside the function body will still
produce ugly errors. Put the constraint on the parameter, not in
the body, so the compiler can reject mismatched callers before
substituting the body.
\end{warning}

\begin{warning}
\textbf{Constrained overloads need partial ordering.} If you have
two function templates with overlapping but distinct concepts,
the compiler picks the most-constrained match. If neither
subsumes the other, the call is ambiguous --- and the error is
not friendly. Compose carefully.
\end{warning}

\begin{warning}
\textbf{Concepts and forward declarations.} You cannot constrain
a forward-declared function template with a concept that
references the function's parameter types if those types depend
on the function itself. The cycle is real; break it by
introducing a helper type.
\end{warning}

\begin{warning}
\textbf{Slow compile times.} Heavy concept use slows the
front-end. \maya{} keeps concepts shallow (a handful of
requirements each) and uses them sparingly outside the public
API. If your build is slow, look for nested concepts more than
two levels deep.
\end{warning}

\section{Concepts as documentation}

The deepest payoff of concepts is not the error message; it is
the documentation. A header file that says

\begin{lstlisting}
template <Program P> int run(Config cfg = {});
\end{lstlisting}

tells you everything you need to know about \code{run}: it takes
a \code{Program} (whose contract is defined elsewhere) and a
config. If you want to know what \code{Program} requires, you go
to its definition --- one place, all the requirements listed.

Compare with the pre-concept C++ idiom:

\begin{lstlisting}
template <class P>  // P must have init(), update(), view(), subscribe()
int run(Config cfg = {});
\end{lstlisting}

The comment is wishful. Nothing enforces it. With concepts, the
comment becomes machine-checked.

\begin{idea}
A concept-driven header is a contract: the compiler is the
referee. When you read \maya{}'s headers, treat the concept
constraints as the most reliable documentation.
\end{idea}

\section*{Workshop}

\begin{enumerate}[leftmargin=*]
  \item Write a concept \code{Printable<T>} that requires
    \code{std::cout << t} to be valid. Use it on a function
    \code{void log(T)}.
  \item Read
    \file{maya/include/maya/core/concepts.hpp} and identify
    every concept declared. Note which other headers they show
    up in.
  \item Try compiling a \code{Counter} app that omits the
    \code{subscribe} method. Read the resulting error and find
    the line that names the failing concept.
\end{enumerate}

\begin{recap}
Concepts are the contracts that make \maya{}'s templates safe and
its error messages readable. Five core concepts shape the public
API. Composing concepts with \code{\&\&} and \code{||} lets you
express precise requirements. Treat them as machine-checked
documentation.
\end{recap}
